\chapter{FUTURE ENHANCEMENTS}
\label{chap:future-work}

\section{Data Resolution and Spatial Processing}
\begin{itemize}
    \item Incorporate daily or pentad (5-day) reanalysis data to capture high-frequency atmospheric triggers, such as westerly wind bursts, which are currently smoothed out by monthly averaging.
    \item Process ORAS5 ocean reanalysis data on its native curvilinear NEMO grid using graph-based or unstructured mesh networks to eliminate bilinear interpolation artifacts in regions with sharp SST gradients (e.g., the eastern Pacific cold tongue).
\end{itemize}

\section{Advanced Spatiotemporal Architectures}
\begin{itemize}
    \item Transition from flattened-tensor tree-based ensembles to end-to-end spatiotemporal deep learning models (e.g., CNN-TCN, Vision Transformers, or Graph Neural Networks) to preserve spatial topology and prevent amplitude underestimation during extreme ENSO events.
    \item Replace discrete Mutual Information Feature Selection (MIFS) with differentiable spatial and channel attention mechanisms to allow the network to dynamically weight feature importance without destroying the 4D tensor structure.
    \item Utilize Graph Neural Networks (GNNs) to process climate data on a true spherical manifold, eliminating spatial distortions introduced by 2D equirectangular grid projections.
\end{itemize}

\section{Probabilistic Forecasting and Uncertainty Quantification}
\begin{itemize}
    \item Shift from deterministic point forecasts to probabilistic frameworks (e.g., Monte Carlo Dropout, Deep Ensembles, or Bayesian Neural Networks) to output predictive distributions.
    \item Generate confidence intervals for Ni\~{n}o 3.4 index and South Asian regional anomaly forecasts to support risk-based decision-making for agriculture and water resource management.
\end{itemize}

\section{Extended Predictive Horizons and Regional Impacts}
\begin{itemize}
    \item Incorporate slow-varying boundary conditions (e.g., global soil moisture, snow cover, Atlantic SSTs) to extend the skillful deterministic forecast horizon beyond the current 3- to 6-month limit and mitigate the spring predictability barrier.
    \item Expand the South Asian regional impact assessment head to forecast additional critical hydro-meteorological variables, including surface runoff, soil moisture, and the frequency of extreme weather events (e.g., tropical cyclones in the Bay of Bengal).
\end{itemize}